\section{Sentiment im Embedding-Raum}

Dieses Kapitel prüft den beschreibenden Teil von Hypothese H1: Ist Sentiment schon vorhanden, bevor das Modell überhaupt eine \gls{transformer}-Schicht durchlaufen hat? Jedes Wort wird zunächst auf einen 1.024-dimensionalen Vektor abgebildet, das sogenannte Embedding. Wenn Sentiment bereits hier erkennbar ist, muss das Modell es nicht erst mühsam aus dem Kontext erschließen.

Eine Hauptkomponentenanalyse reduziert diese 1.024 Dimensionen auf die zwei Achsen, die am meisten Streuung der 2.254 Hu-\&-Liu-Wörter erfassen. Abbildung~\ref{fig:pca} zeigt, dass sich positive und negative Wörter auf diesen zwei Achsen zwar nicht vollständig trennen, aber bereits eine erkennbare Tendenz zeigen: Positive Wörter liegen im Histogramm eher links, negative eher rechts.

\begin{figure}[htbp]
  \centering
  \includegraphics[width=0.49\textwidth]{hu_liu_opinion_embeddings.png}
  \hfill
  \includegraphics[width=0.49\textwidth]{hu_liu_opinion_histogram.png}
  \caption{Hu-\&-Liu-Wörter im Raum der ersten beiden Hauptkomponenten (links) und ihre Verteilung entlang der ersten Hauptkomponente (rechts).}
  \label{fig:pca}
\end{figure}

Dass die Trennung auf nur zwei Achsen unvollständig bleibt, heißt nicht, dass Sentiment nicht vorhanden ist -- es bedeutet nur, dass die Information über viele der 1.024 Dimensionen verteilt ist. Auch die reine Vektorlänge (die L2-Norm) liefert kaum einen Hinweis: Sie liegt für positive und negative Wörter in jedem thematischen Feld nah beieinander (Abbildung~\ref{fig:norms}), wenn auch mit einem leichten, über fast alle Felder konsistenten Muster, dass negative Wörter wie „sour“ oder „dangerous“ eine etwas größere Norm haben als ihre positiven Gegenstücke wie „delicious“ oder „safe“. Aussagekräftiger als die Länge ist die Richtung der Vektoren.

\begin{figure}[htbp]
  \centering
  \includegraphics[width=0.85\textwidth]{sentiment_embedding_norms_all_fields.png}
  \caption{L2-Normen der Sentiment-Embeddings getrennt nach thematischen Feldern.}
  \label{fig:norms}
\end{figure}

Ein ergänzender Blick auf die nächsten Nachbarn im Embedding-Raum, gemessen über die Cosine-Ähnlichkeit (den Winkel zwischen zwei Vektoren, unabhängig von ihrer Länge), bestätigt das Bild: Die einem Wort am nächsten liegenden anderen Wörter besitzen meist dieselbe Sentiment-Polarität, positive Wörter sind also überwiegend von positiven Nachbarn umgeben und negative von negativen. Sogar einfache Embedding-Arithmetik funktioniert in diese Richtung -- addiert man mehrere positive Wörter und zieht ein negatives ab, liegen die nächsten Nachbarn des Ergebnisvektors überwiegend im positiven Bereich, und umgekehrt.

Das zeigt sich noch deutlicher, wenn man gezielt die eine Richtung betrachtet, die aus der Differenz der Embeddings von „good“ und „bad“ gebildet wird: Projiziert man alle Sentiment-Wörter auf diese eine Achse, ergibt sich eine deutlich klarere Trennung. In allen zehn betrachteten thematischen Feldern -- von „Food and taste“ über „Safety and trust“ bis „Reputation and business outcome“ -- liegen positive Wörter fast durchgängig auf der positiven Seite dieser Achse und negative Wörter auf der negativen Seite, mit kaum Überlappung nahe der Mitte. Stärker konnotierte Wörter wie „excellent“ oder „dangerous“ liegen dabei weiter von der Mitte entfernt als schwächer konnotierte wie „pleasant“ oder „rude“. Die Richtung „good“ minus „bad“ verhält sich damit wie eine generelle Sentiment-Achse, die über verschiedenste Themenfelder hinweg funktioniert, obwohl sie nur aus einem einzigen Wortpaar gebildet wurde. Im Feld „Safety and trust“ etwa liegt „dangerous“ deutlich auf der negativen Seite und „trustworthy“ deutlich auf der positiven, während schwächer konnotierte Wörter wie „secure“ oder „insecure“ näher an der Mitte liegen -- die Achse bildet also nicht nur die Richtung, sondern auch die Stärke des Sentiments ab.

Um diese Beobachtung zu quantifizieren, wird ein lineares Probing durchgeführt: Eine logistische Regression -- ein einfacher linearer Klassifikator -- wird auf den Embeddings trainiert und soll positiv von negativ unterscheiden. Gelingt das gut, muss die Information bereits klar im Vektorraum vorliegen, ohne dass eine komplexe Berechnung nötig wäre. Die 2.254 Wörter werden dazu geschichtet im Verhältnis 80 zu 20 aufgeteilt (Tabelle~\ref{tab:split}).

\begin{table}[htbp]
  \centering
  \caption{Geschichtete Aufteilung der Hu-\&-Liu-Wörter für das lineare Probing.}
  \label{tab:split}
  \begin{tabular}{lrrr}
    \toprule
    Teilmenge & positiv & negativ & gesamt \\
    \midrule
    Training & 692 & 1.111 & 1.803 \\
    Test & 173 & 278 & 451 \\
    \midrule
    gesamt & 865 & 1.389 & 2.254 \\
    \bottomrule
  \end{tabular}
\end{table}

Auf den Testdaten erreicht die Probe 97~Prozent Genauigkeit, bei gleichzeitig hohen Präzisions-, Recall- und F1-Werten für beide Klassen (Abbildung~\ref{fig:lraxis}). Die verbleibenden Fehler betreffen erkennbar mehrdeutige oder kontextabhängige Wörter: „capricious“ und „tout“ werden positiv vorhergesagt, obwohl sie im Lexikon als negativ gelten, und „morality“ liegt fast exakt auf der Entscheidungsgrenze -- passend dazu, dass eine moralische Bewertung ohne Kontext kein eindeutiges Sentiment trägt.

\begin{figure}[htbp]
  \centering
  \includegraphics[width=0.8\textwidth]{logistic_regression_sentiment_axis.png}
  \caption{Durch die lineare Probe gelernte Sentiment-Achse; falsch klassifizierte Wörter sind rot markiert.}
  \label{fig:lraxis}
\end{figure}

Damit lässt sich festhalten: Der Embedding-Raum von Pythia-410M enthält bereits eine linear weitgehend separierbare Sentiment-Struktur, bevor überhaupt eine Transformer-Schicht durchlaufen wurde -- der beschreibende Teil von Hypothese H1 ist gestützt. Offen bleibt, ab welcher Schicht das Modell diese Information für seine tatsächliche Vorhersage nutzt -- das untersucht das nächste Kapitel.
